\documentclass[11pt,a4paper]{article}

\usepackage[utf8]{inputenc}
\usepackage[T1]{fontenc}
\usepackage[margin=2.2cm]{geometry}
\usepackage{booktabs}
\usepackage{graphicx}
\graphicspath{{./}}
\usepackage{amsmath}
\usepackage[hidelinks]{hyperref}
\usepackage{caption}
\usepackage{enumitem}
\usepackage{titling}

\captionsetup{font=small,labelfont=bf}
\setlength{\parskip}{0.3em}
\setlength{\parindent}{0pt}
\setlength{\droptitle}{-2.2cm}
\usepackage{titlesec}
\titlespacing*{\section}{0pt}{1.0em}{0.35em}

\title{\textbf{Has Peruvian farmland shifted from food crops to export tree crops,\\
and does legal land title cause it?}\\[0.4em]
\large Short summary of the project to date}
\author{Daniel Ash}
\date{31 August 2026}

\begin{document}
\maketitle
\vspace{-1.2em}

%======================================================================
\section{Purpose and data}

The project asks two questions. Has farmland in Peru moved from annual crops grown for the domestic
market (rice, maize, beans) to perennial tree crops grown for export (mango, lime, coffee, banana)?
And does secure legal title make that move more likely?

The main source is the archive of Peru's land-titling programme, which records for about one
million parcels the boundary, the crop the farmer declared, and the parcel's legal status. It
records the crop only once, mostly between 1997 and 2006. Three further sources were used: a
cadastral bridge file linking the crop registry to the parcel maps, the parcel maps themselves, and
the 2012 agricultural census, which gives a second independent declaration of what is growing on
some of the same parcels. Satellite imagery is available for every year from 1996. The bridge file
covers only 15 departments and one yields nothing usable, so the project covers \textbf{14
departments}, all coastal or lower-altitude. No labels exist for the sierra or the Amazon.

The original plan was to use the archive as training labels, fit a classifier, apply it to every
year of imagery, and read the change from the series. The classifier works. The year-by-year series
does not.

%======================================================================
\section{What worked and what did not}

\begin{table}[htbp]
\centering
\small
\caption{Each line of work and its outcome.}
\label{tab:scoreboard}
\begin{tabular}{@{}p{4.6cm}p{3.3cm}p{7.4cm}@{}}
\toprule
\textbf{Attempted} & \textbf{Outcome} & \textbf{Reason} \\
\midrule
Identify 12 individual crops from imagery
& Fails
& Annual crops share one seasonal curve at 30\,m resolution \\[0.25em]

Separate perennial, annual and bare ground
& Works, 0.68 macro-F1
& Year-round canopy is visible; crop identity is not \\[0.25em]

Extend from Piura to the whole country
& Works, and is more stable
& 14 departments, 946{,}872 parcels \\[0.25em]
Apply the classifier to every year, 1996--2023
& Fails
& A classifier correct 55--59\,\% of the time cannot support a 25-year series per parcel \\[0.25em]

Aggregate that history into 5-year windows, then correct the drift
& Fails; four corrections closed
& The sensor difference depends on land cover, so no single linear correction removes it \\[0.25em]

Photo-interpret recent high-resolution imagery
& Works, 865 parcels labelled
& The only accuracy measured in the period of interest \\[0.25em]

Compare the 1999 declaration with the 2012 census
& Works, 63{,}766 parcels
& Two declarations of the same land, no imagery and no classifier \\[0.25em]

Estimate the causal effect of legal title
& A tight null
& $-0.11$ points, 95\,\% interval $[-1.26, +1.04]$ \\
\bottomrule
\end{tabular}
\end{table}

The classifier separates three land states rather than individual crops: perennial, annual, and
pasture or fallow. That is also the distinction the research question needs, since perennial
approximates export and annual approximates domestic. Gradient-boosted trees (LightGBM) reading
about 135 summary features per parcel-year were selected over a neural network and a hand-written
rule baseline.

%======================================================================
\section{Main results}

\textbf{1. The perennial share rose by about ten percentage points.} Between the titling
declaration (median year 1999) and the 2012 census, the share of parcels declaring a perennial crop
rose from 16.5\,\% to 26.4\,\% nationally, and from 21.0\,\% to 33.5\,\% by land area. The flow is
asymmetric: 7{,}158 parcels moved from annual to perennial against 1{,}832 in the reverse
direction. Conversion happened on larger parcels, median 0.85\,ha against 0.38\,ha for parcels that
stayed annual. Three separate measurements agree on this rise within their uncertainty.

\textbf{2. Legal title does not account for it.} Titled parcels rose 9.2 points and untitled
parcels 11.3 points, so titled parcels shifted slightly less. Measured by area the difference
disappears ($+10.7$ against $+11.0$), which suggests it reflects parcel size rather than tenure. A
difference-in-differences estimate on 14{,}625 parcels and 15 years of imagery returns $-0.11$
percentage points, 95\,\% interval $[-1.26, +1.04]$, and the interval contains zero under every
assumption tested. This is a bounded null rather than proof of no effect: legal status is last
observed in about 2011--12 while the outcome runs to 2023, and part of the comparison group was
probably titled later without being recorded, which pushes any real effect toward zero.

\begin{figure}[htbp]
\centering
\includegraphics[width=\textwidth]{shift_and_tenure.png}
\caption{Left: every source agrees that the perennial share rose, and the two declared-crop
sources agree closely on the size of the rise. The imagery estimates are consistent with them and
too imprecise to add anything. Right: no source finds that legal title accounts for the rise.}
\label{fig:forest}
\end{figure}

\textbf{3. Most land declared as annual cropland in 1999 is not currently cropped.} Read directly
from 865 photo-interpreted parcels, 57.9\,\% of parcels that declared an annual crop are farmable
ground with nothing growing on them, and only 2.9\,\% read as perennial. This is a national average
and does not describe Piura, where 86.7\,\% of such parcels still carry annual crops.
Figure~\ref{fig:transition} shows the full cross-tabulation.

\begin{figure}[htbp]
\centering
\includegraphics[width=0.92\textwidth]{transition_matrix.png}
\caption{Land use declared to the titling programme in 1996--2006 against land use observed on
2019 or later imagery, on 865 hand-labelled parcels, weighted to the national population. Rows sum
to 100\,\%. No classifier is involved. The 95\,\% interval on the 2.9\,\% cell runs from 0 to
5.9.}
\label{fig:transition}
\end{figure}

%======================================================================
\section{Why the year-by-year series failed, and what replaced it}

Two checks were set before the attempt: predictions should degrade gradually when applied to years
far from the training years, and a parcel's predicted history should not alternate implausibly. The
second failed in every version, with 52 to 98\,\% of parcels changing class implausibly often
against a criterion of 15\,\%. Three explanations were ruled out by measurement: prediction quality
does not fall with distance in time, satellite coverage exceeds 93\,\% in all 25 years, and the
failure worsens when the architecture changes. What remains is arithmetic. Errors in different
years are close to independent, so 25 predictions from a classifier correct 55--59\,\% of the time
are dominated by noise. No per-parcel trajectory or area trend was produced from this panel. Two
further obstacles: the 1997--98 El Ni\~no flood makes the baseline years unusable, since perennial
recall on 1998 imagery is 0.000, and five-year aggregation fixes the instability but not the
estimate, because parcels already perennial drift downward faster than the group of interest rises.
The change was therefore measured directly instead, using the 2012 census and photo-interpretation
of recent imagery, neither of which needs a working annual series.

%======================================================================
\section{Method finding}

Cross-validation that holds out nearby land still overstates how well a model performs in a region
it has never seen. Only holding out a whole department reveals the gap, which is about 0.09
macro-F1. This caught three things. Parcel latitude helped by 0.048 macro-F1 on held-out blocks and
hurt by 0.060 when a whole department was held out, so it was acting as a lookup table for
departments already memorised. The neural network won cross-validation in all eight arms and lost
the held-out-department test in all four. And any time-invariant input creates false stability in a
year-by-year model. A second rule follows: score a new input against a same-shaped substitute
rather than against nothing. Climate covariates passed this test for LightGBM and failed it for the
neural network.

%======================================================================
\section{Limitations}

\begin{enumerate}[leftmargin=1.5em,itemsep=1pt]
\item 14 coastal and lower-altitude departments only. No sierra or Amazon.
\item The labels are declarations, not observations, and declaration dates are unreliable.
\item Most accuracy figures come from 1997--2006. Only the recent-imagery figures fall in the
period of interest, and they rest on 865 labels from a single labeller whose consistency has not
been measured.
\item The census link is by farmer name, so it identifies the parcel only approximately and
favours larger, better-documented parcels. Reweighting is applied.
\item One ambiguous class carries a third of the perennial sample. Whether abandoned woody
vegetation counts as perennial moves the 2025 estimates by 8 points on the level and 14 points on
the tenure split, both larger than the effect being measured.
\item Legal status is last observed in 2011--12 while the outcome runs to 2023.
\end{enumerate}

%======================================================================
\section{Next steps}

The main opportunity is to add a properly measured 2025 observation to the 1999 and 2012 points, so
the change can be traced across three dates on the same parcels and split by tenure.
Figure~\ref{fig:tenure} shows the current state of that comparison: the 2025 point is too
imprecise to say anything, and its position depends on how one ambiguous class is treated.
Sentinel-2 gives 19 to 113 clear observations per parcel-year against 13 to 24 for Landsat, so a
single-year 2025 estimate avoids both failures of the year-by-year panel.

\begin{figure}[htbp]
\centering
\includegraphics[width=0.76\textwidth]{tenure_over_time.png}
\caption{The three observations of the same quantity, by tenure at titling. The solid lines are the
census panel. The 2025 points come from a different frame with a different restriction, so each is
drawn from its own baseline rather than continued from the census line. The two imagery readings
differ only in whether abandoned woody vegetation counts as perennial, and that choice moves the
2025 point by more than the effect being measured.}
\label{fig:tenure}
\end{figure}

\begin{enumerate}[leftmargin=1.5em,itemsep=3pt]
\item \textbf{Label more parcels on Sentinel-2 and recent high-resolution imagery.} The campaign
has 865 usable labels and the 2025 intervals are near $\pm$17 points. Keep the existing five
categories and the per-department minimum so a department can still be held out.

\item \textbf{Measure labelling noise first.} Re-label one hundred parcels after an interval.
Report consistency separately for the abandoned orchard against riverside scrub boundary, which is
the hardest call and carries a third of the perennial sample.

\item \textbf{Train a new model on the full 2025 Esri and Sentinel-2 data.} Fit on the new labels
together with the existing ones. Report the held-out-department score as the headline, not the
cross-validation score.

\item \textbf{Run inference on a national sample of parcels for 2025.} Draw from the eligible
universe of 614{,}876 parcels, stratified by department, and reweight to the national population.
The output is the perennial, annual and other split for 2025.

\item \textbf{Place 1999, 2012 and 2025 on one axis for the same parcels.} Restrict to parcels in
the census link, and report the transitions between dates as well as the level at each.

\item \textbf{Split the change by tenure status.} Use legal status at declaration and the
cadastre's 2011--12 status. Report the split by parcel count and by land area, and report it both
with and without abandoned woody vegetation counted as perennial. Later titling in the comparison
group is unobserved, so present the result as a bound rather than a causal estimate.
\end{enumerate}

\end{document}
